%-----------------------------------------------------------------------------------------------
\makeatletter
\immediate\write18{datelog > \jobname.info} % site script for $(date '+%Y-%m-%d %Hh%Mm%Ss')
\makeatother
%-----------------------------------------------------------------------------------------------
%-----------------------------------------------------------------------------------------------
\usetheme{Copenhagen}
\usepackage{beamercolorthemeCNThermSci2}
\usefonttheme{serif}
%-----------------------------------------------------------------------------------------------

%-----------------------------------------------------------------------------------------------
%-----------------------------------------------------------------------------------------------
\usetheme{Copenhagen}
\usepackage{beamercolorthemeUTF2}
\usefonttheme{serif}
%-----------------------------------------------------------------------------------------------
\usepackage[utf8]{inputenc}
\usepackage[greek,french,english,brazil]{babel} % last becomes the active one
\usepackage{pslatex}
\usepackage{amssymb,amsmath}
\usepackage{soul}
\usepackage[squaren,Gray,cdot]{SIunits}
\usepackage[nice]{nicefrac}
\usepackage{tikz}
\usepackage{amscd}
\usepackage{stmaryrd}
\usepackage{scalerel}
\usepackage{xspace}
%-----------------------------------------------------------------------------------------------


%-----------------------------------------------------------------------------------------------
%-----------------------------------------------------------------------------------------------
% Mathematical
%-----------------------------------------------------------------------------------------------
\newcommand{\vet}[1]{\underline{{#1}}}
\newcommand{\mat}[1]{\underline{\underline{{#1}}}}
\newcommand{\cub}[1]{\underline{\underline{\underline{{#1}}}}}
\newcommand{\eqdef}{{\ensuremath\stackrel{\text{\tiny def}}{=}}}
%-----------------------------------------------------------------------------------------------
% Linguistic
%-----------------------------------------------------------------------------------------------
\newcommand{\GRtxt}[1]{\begin{otherlanguage}{greek}{{#1}}\end{otherlanguage}}
\newcommand{\FRtxt}[1]{\begin{otherlanguage}{french}{{#1}}\end{otherlanguage}}
%-----------------------------------------------------------------------------------------------
% Presentation
%-----------------------------------------------------------------------------------------------
\newcommand{\BkgImgH}[1]{% Places an image centered on the slide background filling the height
    \usebackgroundtemplate{\parbox{\paperwidth}{%
        \vspace*{1sp}\centering\includegraphics[height=\paperheight]{{#1}}
}}}
\newcommand{\BkgImgW}[1]{% Places an image centered on the slide background filling the width
    \usebackgroundtemplate{\parbox{\paperwidth}{%
        \vspace*{1sp}\centering\includegraphics[width=\paperwidth]{{#1}}
}}}
\newcommand{\ArtEndH}[3]{% Transitions to plain image (last) slide: #1:prefix #2,#3:extensions
    \BkgImgH{root/../art/#1.#2}
    \frame<handout:0>[plain]{%
        \transdissolve\vspace*{72mm}\color{white}\scriptsize\bf\input{root/../art/#1.#3}}
    \usebackgroundtemplate{\mbox{~}}
}
\newcommand{\ArtEndW}[3]{% Transitions to plain image (last) slide: #1:prefix #2,#3:extensions
    \BkgImgW{root/../art/#1.#2}
    \frame<handout:0>[plain]{%
        \transdissolve\vspace*{72mm}\color{white}\scriptsize\bf\input{root/../art/#1.#3}}
    \usebackgroundtemplate{\mbox{~}}
}
\newcommand{\ImgColW}[3]{% Inserts a full-width image in a column
    \includegraphics[width=\columnwidth]{root/../art/#1.#2}\\[-0.5\baselineskip]
    \parbox{\columnwidth}{\tiny\hfill\scalebox{0.85}{\input{root/../art/#1.#3}}}
}
\newcommand{\txtpic}[1]{%
    \fcolorbox{lightgray}{white!90!black}{{#1}} 
}
%-----------------------------------------------------------------------------------------------


%-----------------------------------------------------------------------------------------------
\title{C.02.01.A1 -- Modelo de Mistura Reativa Ideal}
\subtitle{Aplicação em FTAF -- Finite Time Air-Fuel Otto Engine Model}
\author{Prof.~C.~Naaktgeboren, PhD}
\date{{\scriptsize\tt%
    \includegraphics[height=6.0mm]{root/00-res/cc/by-nc-nd-88x31.pdf}\\[\smallskipamount]
    https://github.com/CNThermSci/ApplThermSci\\
    Compiled on \input{\jobname.info}
}}
%-----------------------------------------------------------------------------------------------
\begin{document}
%-----------------------------------------------------------------------------------------------
\logo{%
    \parbox{158mm}{% There's a 1mm gap on each side of the 160mm x 90mm slide logo line
        \mode<beamer>{
            \includegraphics[height=6.0mm]{root/00-res/UTFPR/UTFPR-logo-D.pdf}\hfill%
            \includegraphics[height=9.0mm]{root/00-res/logo/CNThermSci-logo-A.pdf}%
        }
        \mode<handout>{
            \includegraphics[height=6.0mm]{root/00-res/UTFPR/UTFPR-logo-W.pdf}\hfill%
            \includegraphics[height=9.0mm]{root/00-res/logo/CNThermSci-logo-W.pdf}%
        }
    }
} % The (delineated, alpha), or washed-out logos
%-----------------------------------------------------------------------------------------------
\frame{\titlepage}
%-----------------------------------------------------------------------------------------------

%-----------------------------------------------------------------------------------------------
\section{Modelo de Mistura Reativa Ideal}
%-----------------------------------------------------------------------------------------------

    % !j 96 -i8
    %-------------------------------------------------------------------------------------------
    \begin{frame}{Apresentação: Frações de Massa e Molares:}\vspace*{-2em}
        %-----
        \begin{itemize}
            \item<1-> Mistura \alert{$m$} com \alert{$p$} componentes indexados por \alert{$k$};
            \item<2-> Caracterizada pelas frações mássicas, \alert{$\mathrm{mf}_k$}, e frações
                molares, \alert{$\mathrm{y}_k$}:
        \end{itemize}
        \vspace*{0.8em}\begin{equation*}
            \uncover<3->{
                \alert{\mathrm{mf}_k = \frac{m_k}{m_m}},\qquad\mbox{e}\qquad
                \alert{\mathrm{y}_k  = \frac{n_k}{n_m}},
            }
        \end{equation*}
        %-----
        \begin{itemize}
            \item<4-> Massa da mistura, \alert{$m_m$}, e sua quantidade química, \alert{$n_m$}:
        \end{itemize}
        \vspace*{0.8em}\begin{equation*}
            \uncover<5->{
                \alert{m_m = \sum_{k=1}^p m_k},\qquad\mbox{e}\qquad
                \alert{n_m = \sum_{k=1}^p n_k}.
            }
        \end{equation*}
    \end{frame}
    %-------------------------------------------------------------------------------------------

    % !j 96 -i8
    %-------------------------------------------------------------------------------------------
    \begin{frame}{Massa Molecular e Constante de Gás Aparentes:}\vspace*{-2em}
        \begin{itemize}
            \item<1-> Massa molecular aparente, \alert{$M_m$}: média ponderada pelas
                \alert{frações molares}:
        \end{itemize}
        \vspace*{0.8em}\begin{equation*}
            \uncover<2->{
                \alert{M_m} = \frac{m_m}{n_m} = \alert{\sum_{k=1}^p \mathrm{y}_k M_k},
            }
        \end{equation*}
        \begin{itemize}
            \item<3-> Define a constante de gás aparente, \alert{$R_m$}, junto com $\bar{R}$:
        \end{itemize}
        \vspace*{0.8em}\begin{equation*}
            \uncover<4->{
                \alert{R_m = \frac{\bar{R}}{M_m}}.
            }
        \end{equation*}
    \end{frame}
    %-------------------------------------------------------------------------------------------

    % !j 96 -i8
    %-------------------------------------------------------------------------------------------
    \begin{frame}{Massa Molecular e Constante de Gás Aparentes:}\vspace*{-2em}
        \begin{itemize}
            \item<1-> Mistura de comportamento \alert{$P-T-v\;$ ideal}:
        \end{itemize}
        \vspace*{0.8em}\begin{equation*}
            \uncover<2->{
                \alert{P_mV_m} =
                    \alert{n_m\bar{R}T_m = m_mR_mT_m}.
            }
        \end{equation*}
        \begin{itemize}
            \item<3-> Modelo de calor específico em série truncada:
        \end{itemize}
        \vspace*{0.8em}\begin{equation*}
            \uncover<4->{
                \alert{P_mV_m} =
                    \alert{n_m\bar{R}T_m = m_mR_mT_m}.
            }
        \end{equation*}
    \end{frame}
    %-------------------------------------------------------------------------------------------

    % Finishes with stunning image, with credit
    \ArtEndW{pexels-eberhard-grossgasteiger-542608}{jpg}{txt}

%-----------------------------------------------------------------------------------------------
\end{document}
%-----------------------------------------------------------------------------------------------

